\documentclass[aps,preprint]{revtex4}%
\usepackage{amsfonts}
\usepackage{amsmath}
\usepackage{amssymb}
\usepackage{graphicx}%
\setcounter{MaxMatrixCols}{30}
%TCIDATA{OutputFilter=latex2.dll}
%TCIDATA{Version=5.50.0.2960}
%TCIDATA{CSTFile=revtex4.cst}
%TCIDATA{Created=Saturday, November 07, 2009 09:22:31}
%TCIDATA{LastRevised=Tuesday, November 17, 2009 13:37:02}
%TCIDATA{<META NAME="GraphicsSave" CONTENT="32">}
%TCIDATA{<META NAME="SaveForMode" CONTENT="1">}
%TCIDATA{BibliographyScheme=Manual}
%TCIDATA{<META NAME="DocumentShell" CONTENT="Articles\SW\REVTeX 4">}
%BeginMSIPreambleData
\providecommand{\U}[1]{\protect\rule{.1in}{.1in}}
%EndMSIPreambleData
\newtheorem{theorem}{Theorem}
\newtheorem{acknowledgement}[theorem]{Acknowledgement}
\newtheorem{algorithm}[theorem]{Algorithm}
\newtheorem{axiom}[theorem]{Axiom}
\newtheorem{claim}[theorem]{Claim}
\newtheorem{conclusion}[theorem]{Conclusion}
\newtheorem{condition}[theorem]{Condition}
\newtheorem{conjecture}[theorem]{Conjecture}
\newtheorem{corollary}[theorem]{Corollary}
\newtheorem{criterion}[theorem]{Criterion}
\newtheorem{definition}[theorem]{Definition}
\newtheorem{example}[theorem]{Example}
\newtheorem{exercise}[theorem]{Exercise}
\newtheorem{lemma}[theorem]{Lemma}
\newtheorem{notation}[theorem]{Notation}
\newtheorem{problem}[theorem]{Problem}
\newtheorem{proposition}[theorem]{Proposition}
\newtheorem{remark}[theorem]{Remark}
\newtheorem{solution}[theorem]{Solution}
\newtheorem{summary}[theorem]{Summary}
\newenvironment{proof}[1][Proof]{\noindent\textbf{#1.} }{\ \rule{0.5em}{0.5em}}
\begin{document}
\preprint{HEP/123-qed}
\title[Short title for running header]{Long title that appears on the title page}
\author{First Author}
\affiliation{My Institution}
\author{Second Author}
\affiliation{My Institution}
\author{Third Author}
\affiliation{Other Institution}
\keywords{one two three}
\pacs{PACS number}

\begin{abstract}
Shell document for REV\TeX{} 4.

\end{abstract}
\volumeyear{year}
\volumenumber{number}
\issuenumber{number}
\eid{identifier}
\date[Date text]{date}
\received[Received text]{date}

\revised[Revised text]{date}

\accepted[Accepted text]{date}

\published[Published text]{date}

\startpage{101}
\endpage{102}
\maketitle
\tableofcontents


Consider%

\begin{align}
&  \left(  Y\left\vert \left[  \mathfrak{f}A\right]  _{+}\right.  \right)
=\left(  Y\left\vert \left[  A\mathfrak{f}\right]  _{+}\right.  \right)
=\int_{\boldsymbol{\Delta}_{A}}f\left(  \boldsymbol{r}\right)  \left(
Y\left\vert \left[  \Phi\left(  \boldsymbol{r}\right)  ^{\dagger}\Phi\left(
\boldsymbol{r}\right)  A\right]  _{+}\right.  \right)  d\boldsymbol{r}\\
&  =\int_{\boldsymbol{\Delta}_{A}}f\left(  \boldsymbol{r}\right)
\operatorname{Tr}\left\{  YA\Phi\left(  \boldsymbol{r}\right)  ^{\dagger}%
\Phi\left(  \boldsymbol{r}\right)  \right\}  d\boldsymbol{r}=0;\,\forall
Y\in\mathcal{B}_{2sa}\left(  \mathcal{H}^{N}\right)  ,
\end{align}
where
\begin{equation}
\mathfrak{f=}\int_{\boldsymbol{\Delta}_{A}}f\left(  \boldsymbol{r}\right)
\left\vert \Phi\left(  \boldsymbol{r}\right)  ^{\dagger}\Phi\left(
\boldsymbol{r}\right)  \right)  d\boldsymbol{r}\in\mathcal{B}_{sa}\left(
\mathcal{H}^{N}\right)  .
\end{equation}%
\begin{equation}
\left(  \mathfrak{f}\left\vert B\right.  \right)  =\operatorname{Tr}\left\{
\mathfrak{f}B\right\}  =\int_{\boldsymbol{\Delta}_{A}}f\left(  \boldsymbol{r}%
\right)  \operatorname{Tr}\left\{  \Phi\left(  \boldsymbol{r}\right)
^{\dagger}\Phi\left(  \boldsymbol{r}\right)  B\right\}  dr<\infty\text{. }%
\end{equation}
One can consider $B\in\mathcal{B}_{1sa}\left(  \mathcal{H}^{N}\right)  $ then
$f\left(  \boldsymbol{r}\right)  \in L_{1}\left(  \boldsymbol{\Delta}%
_{A}\right)  $ and $\mathfrak{f}\in\mathcal{B}_{sa}\left(  \mathcal{H}%
^{N}\right)  .$ If $\ f$ is more general then $B$ must belong to a smaller
space i.e. test functions etc.. As $A\in\mathcal{B}_{2sa}\left(
\mathcal{H}^{N}\right)  $ then $A\mathfrak{f}\in\mathcal{B}_{2}\left(
\mathcal{H}^{N}\right)  $ and $\left[  A\mathfrak{f}\right]  _{+}%
\in\mathcal{B}_{2sa}\left(  \mathcal{H}^{N}\right)  $
\begin{align*}
\operatorname{Tr}\left\{  YA\mathfrak{f}\right\}   &  =\left\langle
Y\left\vert A\mathfrak{f}\right.  \right\rangle _{\mathcal{B}_{2}\left(
\mathcal{H}^{N}\right)  }=\frac{1}{2}\operatorname{Tr}\left\{  YA\mathfrak{f}%
+Y\mathfrak{f}A\right\}  =\operatorname{Tr}\left\{  Y\left[  A\mathfrak{f}%
\right]  _{+}\right\}  \\
&  =\left(  \mathfrak{f}\left\vert \widehat{A}\right.  Y\right)  =\left(
Y\left\vert \widehat{A}\right.  \mathfrak{f}\right)  =\left(  Y\left\vert
\left[  A\mathfrak{f}\right]  _{+}\right.  \right)  ,Y\in\mathcal{B}%
_{2sa}\left(  \mathcal{H}^{N}\right)  .
\end{align*}
Also
\begin{align}
\left[  AY\right]  _{+} &  \in\mathcal{B}_{2sa}\left(  \mathcal{H}^{N}\right)
\\
YA &  \in\mathcal{B}_{2}\left(  \mathcal{H}^{N}\right)
\end{align}
and
\begin{align}
\operatorname{Tr}\left\{  \left[  AY\right]  _{+}\mathfrak{f}\right\}   &
=\operatorname{Tr}\left\{  YA\mathfrak{f}\right\}  =\operatorname{Tr}\left\{
M\mathfrak{f}\right\}  \\
M &  =\left[  AY\right]  _{+}.
\end{align}
Then%
\begin{equation}
\operatorname{Tr}\left\{  M\mathfrak{f}\right\}  =\int_{\boldsymbol{\Delta
}_{A}}f\left(  \boldsymbol{r}\right)  \operatorname{Tr}\left\{  M\Phi\left(
\boldsymbol{r}\right)  ^{\dagger}\Phi\left(  \boldsymbol{r}\right)  \right\}
d\boldsymbol{r=}\int_{\boldsymbol{\Delta}_{A}}f\left(  \boldsymbol{r}\right)
\mathcal{M}\left(  \boldsymbol{r}\right)  d\boldsymbol{r}%
\end{equation}
and
\begin{equation}
f\left(  \boldsymbol{r}\right)  \in L_{1}\left(  \boldsymbol{\Delta}%
_{A}\right)  \text{ if }M\in\mathcal{B}_{2sa}\left(  \mathcal{H}^{N}\right)  .
\end{equation}
Using the generalized basis $\left\{  \left\vert \boldsymbol{r}\right\rangle
\right\}  $ we obtain
\begin{align}
\mathfrak{f}\left\vert \boldsymbol{r}\right\rangle  &  =f\left(
\boldsymbol{r}\right)  \left\vert \boldsymbol{r}\right\rangle \\
\mathfrak{f}\left\vert \boldsymbol{r}_{1}\boldsymbol{r}_{2}\right\rangle  &
=\left\{  f\left(  \boldsymbol{r}_{1}\right)  +f\left(  \boldsymbol{r}%
_{2}\right)  \right\}  \left\vert \boldsymbol{r}_{1}\boldsymbol{r}%
_{2}\right\rangle \\
&  \text{etc...}%
\end{align}
then%
\begin{equation}
\mathfrak{f}=\sum_{1\leq N\leq\infty}\binom{\boldsymbol{r}_{N}}{N}^{-1}%
\int_{\boldsymbol{\Delta}_{A}}\sum_{1\leq j\leq N}f\left(  \boldsymbol{r}%
_{j}\right)  \left\vert \boldsymbol{r}_{1}\cdots\boldsymbol{r}_{N}%
\right\rangle \left\langle \boldsymbol{r}_{1}\cdots\boldsymbol{r}%
_{N}\right\vert d\boldsymbol{r}_{1}\cdots d\boldsymbol{r}_{N}%
\end{equation}
and as $A\in\mathcal{B}_{2sa}\left(  \mathcal{H}^{N}\right)  $%
\begin{equation}
A=\sum_{1\leq j\leq r_{N}}\alpha_{j}\left\vert \Psi_{j}\right\rangle
\left\langle \Psi_{j}\right\vert ;\,\alpha_{j}\in\mathbb{R},\left\vert
\Psi_{j}\right\rangle \in\mathcal{H}^{N},1\leq j\leq r_{N}.
\end{equation}
Then the linear independence condition becomes
\begin{align}
&  \sum_{1\leq j\leq N}\int_{\boldsymbol{\Delta}_{A}}f\left(  \boldsymbol{r}%
_{j}\right)  \left\vert \boldsymbol{r}_{1}\cdots\boldsymbol{r}_{N}%
\right\rangle \left\langle \boldsymbol{r}_{1}\cdots\boldsymbol{r}%
_{N}\right\vert d\boldsymbol{r}_{1}\cdots d\boldsymbol{r}_{N}\sum_{1\leq j\leq
r_{N}}\alpha_{j}\left\vert \Psi_{j}\right\rangle \left\langle \Psi
_{j}\right\vert \\
&  =\sum_{1\leq j\leq r_{N}}\alpha_{j}\sum_{1\leq j\leq N}\int%
_{\boldsymbol{\Delta}_{A}}f\left(  \boldsymbol{r}_{j}\right)  \left\langle
\boldsymbol{r}_{1}\cdots\boldsymbol{r}_{N}\left\vert \Psi_{j}\right.
\right\rangle \left\vert \boldsymbol{r}_{1}\cdots\boldsymbol{r}_{N}%
\right\rangle \left\langle \Psi_{j}\right\vert d\boldsymbol{r}_{1}\cdots
d\boldsymbol{r}_{N}\\
&  =\sum_{1\leq k\leq r_{N}}\alpha_{k}\sum_{1\leq j\leq N}\int%
_{\boldsymbol{\Delta}_{A}}f\left(  \boldsymbol{r}_{j}\right)  \Psi_{k}\left(
\boldsymbol{r}_{1}\cdots\boldsymbol{r}_{N}\right)  \left\vert \boldsymbol{r}%
_{1}\cdots\boldsymbol{r}_{N}\right\rangle \left\langle \Psi_{k}\right\vert
d\boldsymbol{r}_{1}\cdots d\boldsymbol{r}_{N}\\
&  =0,
\end{align}
giving
\begin{align}
&  \alpha_{k}\left\{  \sum_{1\leq j\leq N}f\left(  \boldsymbol{r}_{j}\right)
\right\}  \Psi_{k}\left(  \boldsymbol{r}_{1}\cdots\boldsymbol{r}_{N}\right)
=0\,\forall k\\
&  \Rightarrow\alpha_{k}\Psi_{k}\left(  \boldsymbol{r}_{1}\cdots
\boldsymbol{r}_{N}\right)  =0\text{ }\forall k\text{ when }\sum_{1\leq j\leq
N}f\left(  \boldsymbol{r}_{j}\right)  \neq0
\end{align}
As $\Psi_{k}\left(  \boldsymbol{r}_{1}\cdots\boldsymbol{r}_{N}\right)
\overset{a.e.}{\neq}0$ when $\left\{  \boldsymbol{r}_{1}\cdots\boldsymbol{r}%
_{N}\right\}  \in\boldsymbol{\Delta}_{A}\times\cdots\times\boldsymbol{\Delta
}_{A}$ this produces a contradiction, thus
\begin{equation}
\sum_{1\leq j\leq N}f\left(  \boldsymbol{r}_{j}\right)  \overset{a.e.}{=}0.
\end{equation}


\bigskip The operators $\left\{  \left.  \left\vert \left[  \Phi\left(
\boldsymbol{r}\right)  ^{\dagger}\Phi\left(  \boldsymbol{r}\right)  A\right]
_{+}\right)  \right\vert \mathbf{r\in}\mathbb{R}\right\}  $ can be
orthonormalized
\begin{equation}
\left\vert \left[  \Phi^{\dagger}\Phi A\right]  _{+}\right)  \left(  \left.
\left[  \Phi^{\dagger}\Phi A\right]  _{+}\right\vert \left[  \Phi^{\dagger
}\Phi A\right]  _{+}\right)  ^{-\frac{1}{2}}%
\end{equation}
i.e. %

\begin{align}
&  \int_{\boldsymbol{\Delta}_{A}}\left\vert \left[  \Phi\left(  \boldsymbol{r}%
^{\prime}\right)  ^{\dagger}\Phi\left(  \boldsymbol{r}^{\prime}\right)
A\right]  _{+}\right)  \left(  \left.  \left[  \Phi\left(  \boldsymbol{r}%
^{\prime}\right)  ^{\dagger}\Phi\left(  \left(  \boldsymbol{r}^{\prime
}\right)  \right)  A\right]  _{+}\right\vert \left[  \Phi\left(
\boldsymbol{r}\right)  ^{\dagger}\Phi\left(  \boldsymbol{r}\right)  A\right]
_{+}\right)  ^{-\frac{1}{2}}d\boldsymbol{r}^{\prime}\\
&  =\frac{1}{2}\int_{\boldsymbol{\Delta}_{A}}\left\vert \left[  \Phi\left(
\boldsymbol{r}^{\prime}\right)  ^{\dagger}\Phi\left(  \boldsymbol{r}^{\prime
}\right)  A\right]  _{+}\right)  \left(  \operatorname{Tr}\left\{  \Phi\left(
\boldsymbol{r}^{\prime}\right)  ^{\dagger}\Phi\left(  \boldsymbol{r}^{\prime
}\right)  A\Phi\left(  \boldsymbol{r}\right)  ^{\dagger}\Phi\left(
\boldsymbol{r}\right)  A+\Phi\left(  \boldsymbol{r}^{\prime}\right)
^{\dagger}\Phi\left(  \boldsymbol{r}^{\prime}\right)  \Phi\left(
\boldsymbol{r}\right)  ^{\dagger}\Phi\left(  \boldsymbol{r}\right)
A^{2}\right\}  \right)  ^{-\frac{1}{2}}d\boldsymbol{r}^{\prime}%
\end{align}
produced by the projector
\begin{equation}
\mathcal{P}=\int_{\boldsymbol{\Delta}_{A}}\frac{\left\vert \left[  \Phi\left(
\boldsymbol{r}\right)  ^{\dagger}\Phi\left(  \boldsymbol{r}\right)  A\right]
_{+}\right)  \left(  \left[  \Phi\left(  \boldsymbol{r}\right)  ^{\dagger}%
\Phi\left(  \boldsymbol{r}\right)  A\right]  _{+}\right\vert }{\left(  \left.
\left[  \Phi\left(  \boldsymbol{r}\right)  ^{\dagger}\Phi\left(
\boldsymbol{r}\right)  A\right]  _{+}\right\vert \left[  \Phi\left(
\boldsymbol{r}\right)  ^{\dagger}\Phi\left(  \boldsymbol{r}\right)  A\right]
_{+}\right)  }d\boldsymbol{r}\nonumber
\end{equation}
whose action is given by
\begin{align}
\mathcal{P}\left\vert B\right)   &  =\int_{\boldsymbol{\Delta}_{A}}%
\frac{\left(  \left.  \left[  \Phi\left(  \boldsymbol{r}\right)  ^{\dagger
}\Phi\left(  \boldsymbol{r}\right)  A\right]  _{+}\right\vert B\right)
}{\left(  \left.  \left[  \Phi\left(  \boldsymbol{r}\right)  ^{\dagger}%
\Phi\left(  \boldsymbol{r}\right)  A\right]  _{+}\right\vert \left[
\Phi\left(  \boldsymbol{r}\right)  ^{\dagger}\Phi\left(  \boldsymbol{r}%
\right)  A\right]  _{+}\right)  }\left\vert \left[  \Phi\left(  \boldsymbol{r}%
\right)  ^{\dagger}\Phi\left(  \boldsymbol{r}\right)  A\right]  _{+}\right)
d\boldsymbol{r}\nonumber\\
&  =\int_{\boldsymbol{\Delta}_{A}}\frac{\operatorname{Tr}\left\{  \left[
\Phi\left(  \boldsymbol{r}\right)  ^{\dagger}\Phi\left(  \boldsymbol{r}%
\right)  A\right]  _{+}B\right\}  }{\operatorname{Tr}\left\{  \left[
\Phi\left(  \boldsymbol{r}\right)  ^{\dagger}\Phi\left(  \boldsymbol{r}%
\right)  A\right]  _{+}\left[  \Phi\left(  \boldsymbol{r}\right)  ^{\dagger
}\Phi\left(  \boldsymbol{r}\right)  A\right]  _{+}\right\}  }\left\vert
\left[  \Phi\left(  \boldsymbol{r}\right)  ^{\dagger}\Phi\left(
\boldsymbol{r}\right)  A\right]  _{+}\right)  d\boldsymbol{r}\nonumber\\
&  =\int_{\boldsymbol{\Delta}_{A}}\operatorname{Tr}\left\{  A\Phi\left(
\boldsymbol{r}\right)  ^{\dagger}\Phi\left(  \boldsymbol{r}\right)  B\right\}
\left\vert \left[  \Phi\left(  \boldsymbol{r}\right)  ^{\dagger}\Phi\left(
\boldsymbol{r}\right)  A\right]  _{+}\right)  \frac{1}{\rho_{A^{2}}\left(
\boldsymbol{r}\right)  }d\boldsymbol{r}.
\end{align}



\end{document}